\documentclass[11pt,a4paper]{article}
\usepackage[T1]{fontenc}
\usepackage[utf8]{inputenc}
\usepackage{lmodern,amsmath,amssymb}
\usepackage[margin=25mm]{geometry}
\usepackage{longtable,booktabs,array,calc}
\usepackage[hidelinks]{hyperref}
\hypersetup{pdftitle={Galileo's comparison as two-path interference},pdfauthor={navstokgap research working notes}}
\newcommand{\tightlist}{\setlength{\itemsep}{0pt}\setlength{\parskip}{0pt}}
\setlength{\emergencystretch}{3em}
\setcounter{secnumdepth}{0}
\begin{document}
\hypertarget{galileos-comparison-as-two-path-interference}{%
\section{Galileo's comparison as two-path
interference}\label{galileos-comparison-as-two-path-interference}}

\textbf{After refereeing (Fable, 2026-09-27).} Lemma 1, the conversions
\(K_\tau=F^2\tau^3/(24m)=(F/4v)A=\tau\Delta E/12\), the error law and
the polygon-phase match, and Theorem 3(a),(c) are accepted. Theorem 3(b)
holds for the relative-averaging prescription stated there (see after
(12)). ``Two-path interference'' here means a controlled-unitary
comparison, and the \(n\) passes are product copies; Proposition 2 is
not used by Theorem 3.

\textbf{Result, 2026-09-27 (round 3).} A coherent, endpoint-matched
comparison of Newton's impulsive chord with Galileo's constant-force
parabola has action difference

\[
K_\tau:=\Delta S=\frac{F^2\tau^3}{24m}
=\frac{F}{4v}A_{\rm inertial,fall}=\frac{\tau\Delta E}{12}.
\]

At supplied action resolution \(\varepsilon>0\), its optimal equal-prior
error against the zero-phase reference, using \(n\) independent
controlled passes and an arbitrary joint measurement, is exactly

\[
p_{\rm err}^{(n)}=\frac12\left(1-
\sqrt{1-\cos^{2n}\!\left(\frac{K_\tau}{2\varepsilon}\right)}\right).
\]

For fixed \(n,\varepsilon\), refinement sends this error to \(1/2\). At
fixed \(\tau>0\), the coherent error oscillates as
\(\varepsilon\downarrow0\) and has no limit. A specified averaging of
the inverse resolution gives a genuine noncommutation theorem for
coherence; a specified tunable comparison gives one for
distinguishability (Theorem 3). These prescriptions make precise the
extra operation needed in the proposed classical-first limit.

The resulting action threshold is conditional on the phase rule, the
Born rule and the resources of the comparison. It re-expresses the
polygon-lift phase bound through the halved functional. Theorems C and E
of the fifth-postulate note constrain different records with the same
supplied quantum action constant. Positivity of that constant remains an
independent physical obligation. All results below are proved under
their displayed hypotheses; the scale-selection step remains open.

\hypertarget{the-galileo-area-and-the-two-histories}{%
\subsection{1. The Galileo area and the two
histories}\label{the-galileo-area-and-the-two-histories}}

Let \(m,F,v,\tau>0\), with downward coordinate \(y\) and common
horizontal motion \(x=vt\). The initial tangent and falling curve are

\[
y_I(t)=0,\qquad y_F(t)=\frac{Ft^2}{2m},\qquad
A_{\rm inertial,fall}=\frac{vF\tau^3}{6m},\qquad
\Delta E=\frac{F^2\tau^2}{2m}.
\tag{1}
\]

Thus \((3F/v)A_{\rm inertial,fall}=\tau\Delta E\). The common-endpoint
chord is \(y_C(t)=F\tau t/(2m)\). For the transverse Lagrangian
\(L=m\dot y^2/2+Fy\), direct integration gives

\[
S[y_I]=0,\qquad S[y_F]=\frac{F^2\tau^3}{3m},\qquad
S[y_C]=\frac{3F^2\tau^3}{8m},\qquad S[y_C]-S[y_F]=K_\tau.
\tag{2}
\]

The horizontal action cancels in every difference. Equation (2) also
exposes the endpoint issue: the bare tangent and parabola finish at
different heights. Adding \(d\chi(y,t)/dt\) to \(L\) changes their
action difference by their unequal final \(\chi\) values. Their
measurable relative phase therefore requires endpoint matching and
specified control phases. The conversion coefficient for \(K_\tau\) is
\(F/(4v)\); the area conversion \((F/v)A_{\rm inertial,fall}\) equals
\(4K_\tau\).

Here is the exact matching used throughout. Supply canonical quantum
kinematics \([y,p]=i\varepsilon\) in the Schrödinger representation, and
set

\[
U_\tau=\exp\left[-\frac{i\tau}{\varepsilon}
\left(\frac{p^2}{2m}-Fy\right)\right],\qquad
Q_\tau=e^{iF\tau y/(2\varepsilon)}
e^{-i\tau p^2/(2m\varepsilon)}e^{iF\tau y/(2\varepsilon)}.
\]

\textbf{Lemma 1 (matched comparison).} With these scalar-potential
conventions,

\[Q_\tau=e^{iK_\tau/\varepsilon}U_\tau.\tag{3}\]

\emph{Proof.} The kernel of \(Q_\tau\) has the free normalization and
phase \(L_\tau(x,y)=m(y-x)^2/(2\tau)+F\tau(x+y)/2\). The classical
solution joining \(x\) to \(y\) is
\(q_*(t)=x+(y-x)t/\tau+F(t^2-\tau t)/(2m)\). Substitution into the
action gives \(L_\tau-K_\tau\). The fluctuation action about \(q_*\) is
exactly \(\int m\dot\xi^2/2\), with zero endpoint fluctuations, hence
the same free Gaussian normalization. This proves (3), equivalently the
exact quadratic kernel identity in the
\href{refinement-composition-and-limit.md}{insertion note, §§2--3}
(passage). \(\square\)

The impulses in \(Q_\tau\) reproduce the chord and match the parabola's
final momentum after the last half-kick. Equation (3) holds for every
initial body state. A branch-dependent scalar potential would add its
own phase, so its absence or calibration is part of the experiment. The
\href{polygon-lift-phase.md}{polygon-lift note, Theorem 1} (full-read)
proves the corresponding identity for every partition.

For a literal tangent arm, the explicit reunion in
\href{newton-insertion-action.md}{the Galileo note, §§3--4} (full-read)
uses signal forces \(F,-F,F\) over durations \(\tau,2\tau,\tau\) against
a free reference. Its closed action difference is
\(\Delta S_{\rm loop}=-2F^2\tau^3/(3m)=-16K_\tau\). Every phase formula
below applies with that signed difference when that four-cell experiment
is chosen. The coefficient belongs to the reunion protocol; the original
area (1) remains the geometric anchor.

For fixed endpoints the constant-force action has one stationary path.
The chord is an off-shell path for the continuous-force action and a
classical history for the impulsive control. Thus (3) compares two
controlled histories, each with an exact quadratic path integral.
Treating the tangent and parabola as two stationary points of one
fixed-endpoint constant-force functional would change the variational
problem.

\hypertarget{what-the-halved-functional-measures}{%
\subsection{2. What the halved functional
measures}\label{what-the-halved-functional-measures}}

Use \(\varepsilon\) for action resolution and \(\eta\in[0,1/2)\) for a
target decision error. Identification with quantum mechanics means
\(\varepsilon=|\hbar|\), with ordinary Planck constant
\(h_P=2\pi|\hbar|\).

\textbf{Proposition 2 (two contributions and their remainder).} Suppose
a halved functional, with its normalization and observable fixed, has
two retained contributions

\[
\mathcal A_\varepsilon
=a_0e^{iS_0/\varepsilon+i\mu_0}
+a_1e^{iS_1/\varepsilon+i\mu_1}+R_\varepsilon,
\quad a_j\ge0,\quad |R_\varepsilon|\le r.
\]

Put \(\phi=(S_1-S_0)/\varepsilon\) and \(\beta=\mu_1-\mu_0\). Then

\[
|\mathcal A_\varepsilon|^2
=a_0^2+a_1^2+2a_0a_1\cos(\phi+\beta)+E,
\qquad |E|\le2(a_0+a_1)r+r^2.\tag{4}
\]

\emph{Proof.} Expand the square, bound the leading amplitude by
\(a_0+a_1\), and bound its cross term with \(R_\varepsilon\) by
\(2(a_0+a_1)r\). \(\square\)

For two nondegenerate stationary points at a fixed partition,
\(a_j=|O(z_j)|/\sqrt{|\det S''(z_j)|}\) and \(\mu_j\) contains the
observable phase and Hessian signature. Stationary phase supplies
\(r\le C_\pi\varepsilon\) for sufficiently small \(\varepsilon\), with
cutoffs and smoothness as in
\href{rivero-1998-conjecture-central-forces.md}{round 1, §1}
(full-read). Equation (4) retains that partition-dependent constant;
uniform control as the partition changes needs a separate estimate. For
the exactly matched quadratic propagators in (3), the relative phase has
zero remainder and zero relative Gaussian signature.

At a fixed phase, a large \(|S_1-S_0|/\varepsilon\) still leaves the
cross term in (4) at full amplitude. Averaging, selecting a branch, or
retaining an orthogonal branch record supplies the missing
diagonalization. Equal-action contributions survive phase averaging, as
round 1's pair of quartic saddles demonstrates.

An operational record requires a preparation and measurement rule in
addition to (4). For example, tag the arms by normalized pointer states
\(|r_0\rangle,|r_1\rangle\), and put \(\gamma=\langle r_0|r_1\rangle\).
Recombining the arms gives the cross term
\(2a_0a_1\operatorname{Re}(\gamma e^{i(\phi+\beta)})\) and visibility
\(V=2a_0a_1|\gamma|/(a_0^2+a_1^2)\). For balanced arms and pure tags,
optimal inference of the arm from the pointer has

\[
D_{\rm path}=\sqrt{1-|\gamma|^2},\qquad
p_{\rm path}=\frac{1-D_{\rm path}}2,\qquad
V^2+D_{\rm path}^2=1.\tag{5}
\]

The eigenvalues of the difference of the two rank-one tag states are
\(\pm\sqrt{1-|\gamma|^2}\), proving (5) by the binary measurement
optimization used below. This is the pure balanced case of established
interferometric duality; see
\href{https://arxiv.org/abs/quant-ph/9908072}{Schwindt, Kwiat and
Englert 1999} (abstract). Path identification depends on the tags, while
phase detection depends on coherent recombination. Orthogonal tags
identify the arm even at \(\Delta S=0\) and remove its interference. A
bare halved functional supplies neither those tags nor their physical
cost.

\hypertarget{exact-phase-record-error-target-and-resources}{%
\subsection{3. Exact phase record, error target and
resources}\label{exact-phase-record-error-target-and-resources}}

Prepare a balanced control qubit and a common arbitrary body density
operator \(\rho\). Under hypothesis \(H_0\), apply \(U_\tau\) in both
arms; under \(H_1\), apply \(U_\tau\) in arm 0 and \(Q_\tau\) in arm 1.
Supply coherent control with calibrated relative scalar phases, no extra
environmental tags, the Born rule, and equal hypothesis priors. Equation
(3) factorizes the output into the same body state under both hypotheses
and the control states

\[
|+\rangle=\frac{|0\rangle+|1\rangle}{\sqrt2},\qquad
|+_\phi\rangle=\frac{|0\rangle+e^{i\phi}|1\rangle}{\sqrt2},
\qquad \phi=K_\tau/\varepsilon.\tag{6}
\]

A recombiner of phase \(\alpha\) has exact port probabilities
\(P_\pm(\alpha|\phi)=[1\pm\cos(\phi-\alpha)]/2\). Scanning \(\alpha\)
reveals the fringe shift; its visibility remains one for every \(\phi\).
The optimal distinction of the two hypotheses is

\[
D_n=\sqrt{1-\cos^{2n}(\phi/2)},\qquad
p_{\rm err}^{(n)}=\frac{1-D_n}{2}.\tag{7}
\]

\emph{Proof.} For normalized vectors with overlap \(c\), the difference
of their rank-one density operators has eigenvalues
\(\pm\sqrt{1-|c|^2}\). For any equal-prior binary test \(0\le M\le I\),
the success probability is \([1+\operatorname{tr}M(\rho_1-\rho_0)]/2\);
the positive eigenspace attains the maximum. Here
\(|\langle+|+_\phi\rangle|=|\cos(\phi/2)|\) and \(n\) independent copies
raise the overlap to its \(n\)th power. This proves (7) including
attainability. \(\square\)

Thus (7) detects a history-dependent phase against a reference. The
uncontrolled body channels of \(Q_\tau\) and \(U_\tau\) coincide for
every \(\rho\), giving error \(1/2\) even with arbitrarily many copies.
The coherence resource in (6) is essential to exposing the scalar.

For the original unclosed free/falling pair, the final displacement and
impulse themselves supply a further signal. A pure body's overlap then
contains the expectation of a Weyl displacement in its prepared state,
so its magnitude depends on that preparation. For each \(\tau>0\) one
can choose a wavefunction with position support narrower than the
nonzero interaction-picture displacement \(F\tau^2/(2m)\); its
translated support is disjoint and the two outputs are orthogonal.
Shrinking that support costs increasing momentum spread. Thus the
phase-only theorem applies to the matched, coherently controlled record
(6), with its specified resources; an unrestricted detector of the raw
endpoint separation requires a separate preparation or disturbance
budget.

Write \(\delta=\operatorname{dist}(\phi,2\pi\mathbb Z)\in[0,\pi]\). For
a fixed positive integer \(n\), (7) reaches error at most \(\eta\) iff

\[
\delta\ge d_{n,\eta},\qquad
d_{n,\eta}=2\arccos\left([4\eta(1-\eta)]^{1/(2n)}\right).
\tag{8}
\]

On the first lobe \(0\le K_\tau/\varepsilon\le\pi\), this is equivalent
to

\[
K_\tau\ge\varepsilon d_{n,\eta},\qquad
\tau\ge\left(\frac{24m\varepsilon d_{n,\eta}}{F^2}\right)^{1/3},
\qquad A_{\rm inertial,fall}\ge\frac{4v\varepsilon}{F}d_{n,\eta}.
\tag{9}
\]

For \(n=1\), \(d_{1,\eta}=2\arcsin(1-2\eta)\), exactly the
\href{polygon-lift-phase.md}{polygon-lift record bound, §4} (passage).
For \(0<\eta<1/2\), \(d_{n,\eta}\to0\) as \(n\to\infty\); for \(\eta=0\)
it equals \(\pi\) at every finite \(n\). With \(k\) coherent passes
replace \(\phi\) by \(k\phi\). These resource choices change the
threshold.

The dependence on errors is explicit. If actual hypothesis states
\(\widetilde\rho_i\) satisfy
\(\tfrac12\|\widetilde\rho_i-\rho_i\|_1\le e_i\), the optimum error
changes by at most \((e_0+e_1)/2\), by the triangle inequality in (7). A
common phase calibration error \(u\) on each of \(n\) independent signal
copies gives distance at most
\(\sqrt{1-\cos^{2n}(u/2)}\le\sqrt n\,|u|/2\), and hence changes the
optimal error by at most \(\sqrt n\,|u|/4\) when the reference is exact.
An action uncertainty \(|\delta S|\) contributes
\(|u|=|\delta S|/\varepsilon\). These bounds state how accurately the
control must be known to use (9).

\hypertarget{the-two-orders-of-limits}{%
\subsection{4. The two orders of
limits}\label{the-two-orders-of-limits}}

Put \(C=F^2/(24m)>0\), so \(K_\tau=C\tau^3\).

\textbf{Theorem 3 (refinement, classical resolution and the missing
prescription).}

\begin{enumerate}
\def\labelenumi{(\alph{enumi})}
\tightlist
\item
  For the fixed coherent protocol (6), with fixed finite \(n\),
\end{enumerate}

\[
D_n\le\min\left(1,\frac{\sqrt n\,C\tau^3}{2\varepsilon}\right),
\qquad \lim_{\tau\downarrow0}p_{\rm err}^{(n)}=\frac12
\quad(\varepsilon>0).\tag{10}
\]

At every fixed \(\tau>0\), \(\lim_{\varepsilon\downarrow0}D_n\) fails to
exist: its liminf is zero and limsup is one. Consequently the coherent
classical-first iterated limit is undefined.

\begin{enumerate}
\def\labelenumi{(\alph{enumi})}
\setcounter{enumi}{1}
\tightlist
\item
  Fix \(b\in(0,1)\) and average the inverse resolution by replacing
  \(\phi\) with \(u\phi\), uniformly over \(u\in[1-b,1+b]\), forgetting
  \(u\). For this explicitly averaged two-arm state, its normalized
  coherence is
\end{enumerate}

\[
\Gamma_b(\tau,\varepsilon)
=\frac1{2b}\int_{1-b}^{1+b}e^{iuK_\tau/\varepsilon}\,du
=e^{i\phi}\operatorname{sinc}(b\phi),\qquad
|\Gamma_b|\le\min\left(1,\frac{\varepsilon}{bK_\tau}\right).
\tag{11}
\]

With \(\operatorname{sinc}(0)=1\), the iterated limits are

\[
\lim_{\varepsilon\downarrow0}\lim_{\tau\downarrow0}\Gamma_b=1,
\qquad
\lim_{\tau\downarrow0}\lim_{\varepsilon\downarrow0}\Gamma_b=0.
\tag{12}
\]

(Refinement after a Fable review, 2026-09-27.) The non-commuting limits
(12) depend on averaging over a \emph{relative} window of the inverse
resolution. With an absolute window of half-width \(w_0\) in
\(1/\varepsilon\) the average is
\(e^{i\phi}\operatorname{sinc}(K_\tau w_0)\), independent of
\(\varepsilon\), and both iterated limits equal one. Theorem 3(b)
therefore states a property of the relative-averaging prescription,
which is the prescription (12) uses.

For fixed weights and signature difference in (4), this prescription
leaves a diagonal-intensity error bounded by
\(2a_0a_1\min(1,\varepsilon/(b|\Delta S|))\) in addition to the averaged
remainder bound of (4), with bound \(2a_0a_1\) when \(\Delta S=0\). If
the stationary-phase remainder obeys \(r\le C_\pi\varepsilon\) uniformly
for small resolution, averaging uses resolutions \(\varepsilon/u\) and
one may take \(r\le C_\pi\varepsilon/(1-b)\).

\begin{enumerate}
\def\labelenumi{(\alph{enumi})}
\setcounter{enumi}{2}
\tightlist
\item
  Alternatively, assume a calibrated comparison can be attenuated to any
  action \(rK_\tau\), \(0\le r\le1\), under both hypotheses, with no
  additional relative phase. For example, use duration \(r^{1/3}\tau\)
  in (6), keeping \(F,m\) fixed. Let \(D_{\rm tun}\) be the best
  single-copy distance over these choices. Then
\end{enumerate}

\[
D_{\rm tun}(\tau,\varepsilon)
=\sin\left(\min\left\{\frac{K_\tau}{2\varepsilon},\frac\pi2\right\}\right),
\tag{13}
\] \[
\lim_{\varepsilon\downarrow0}\lim_{\tau\downarrow0}D_{\rm tun}=0,
\qquad
\lim_{\tau\downarrow0}\lim_{\varepsilon\downarrow0}D_{\rm tun}=1.
\tag{14}
\]

\emph{Proof.} The inequality \(1-(1-x)^n\le nx\) for \(x\in[0,1]\), with
\(x=\sin^2(\phi/2)\), proves (10). For fixed \(K_\tau>0\), the sequences
\(\varepsilon_j=K_\tau/(2\pi j)\) and
\(\varepsilon'_j=K_\tau/((2j+1)\pi)\) give respectively \(D_n=0\) and
\(D_n=1\). Integration gives (11); dominated continuity at \(\phi=0\)
and the displayed reciprocal bound at \(\phi\to\infty\) prove (12). In
(c), maximize \(|\sin(r\phi/2)|\) over \([0,1]\). The first maximum is
reached at \(r=\pi/\phi\) when \(\phi\ge\pi\); otherwise \(r=1\)
maximizes it. This proves (13)--(14). \(\square\)

Part (b) formalizes the loss of off-diagonal terms needed to recover the
two selected critical-point weights. Its averaged balanced state
approaches \(I/2\) classically first, and \(|+\rangle\langle+|\) when
refined first. The trace distance between these two limiting states is
\(1/2\); dephasing by itself supplies no accessible which-path tag. Part
(c) instead supplies an actual discrimination experiment with errors
tending respectively to \(1/2\) and zero. Its extra tunable control is a
physical hypothesis; changing the output recombiner alone would leave
(7) unchanged. A classical branch-selection prescription can also retain
distinct histories at every positive cell length, but their paths
converge together as that length vanishes. A limiting Dirac measure
alone expresses no retained record of the shrinking difference.

These conclusions sharpen the proposed classical-first statement. For
any fixed \(b\) the condition \(b|\Delta S|/\varepsilon\gg1\) controls
the averaged cross term, with the explicit error in (11). For coherent
discrimination the exact condition is (8), with its phase aliases. The
two conditions express different measurements. Along a joint limit the
dimensionless ratio \(C\tau^3/\varepsilon\) governs both: tending to
zero erases the phase signal; a finite ratio gives the corresponding
signal, including the aliases in (8); a diverging ratio requires the
prescription in (b) or (c) for a definite limit.

For a partition of fixed total duration \(T\), the matched polygon
action is \(K_\pi=C\sum_j\tau_j^3\le CT|\pi|^2\). Its coherent distance
from the exact curve is at most \(\sqrt n\,CT|\pi|^2/(2\varepsilon)\).
After \(k\) uniform halvings a cell's action falls by \(8^{-k}\) and the
whole polygon defect by \(4^{-k}\). Hence the whole fixed-duration
comparison converges as well. Resolving ever smaller cells can require
increasing resources: for a fixed small-error target the small-phase
independent-copy scaling is \(n\) of order \((\varepsilon/K_\tau)^2\),
whereas coherent repetition uses \(k\) of order \(\varepsilon/K_\tau\).
Such budgets were held fixed in (10).

\hypertarget{relation-to-the-disturbance-and-concentration-floors}{%
\subsection{5. Relation to the disturbance and concentration
floors}\label{relation-to-the-disturbance-and-concentration-floors}}

The \href{principia-fifth-postulate.md}{fifth-postulate note, Theorems
C, E and §7} (passage) separates observable algebra, states and records.
With \(\varepsilon=|\hbar|\), the three bounds here have the following
scopes.

\begin{longtable}[]{@{}
  >{\raggedright\arraybackslash}p{(\columnwidth - 4\tabcolsep) * \real{0.3333}}
  >{\raggedright\arraybackslash}p{(\columnwidth - 4\tabcolsep) * \real{0.3333}}
  >{\raggedright\arraybackslash}p{(\columnwidth - 4\tabcolsep) * \real{0.3333}}@{}}
\toprule\noalign{}
\begin{minipage}[b]{\linewidth}\raggedright
Bound
\end{minipage} & \begin{minipage}[b]{\linewidth}\raggedright
Quantity constrained
\end{minipage} & \begin{minipage}[b]{\linewidth}\raggedright
Extra hypotheses beyond Newtonian geometry
\end{minipage} \\
\midrule\noalign{}
\endhead
\bottomrule\noalign{}
\endlastfoot
Phase record (9), one copy & \(K_\tau\ge2|\hbar|\arcsin(1-2\eta)\) on
the first lobe & Coherent controlled comparison, calibrated reunion,
Born rule and one pass \\
Theorem C(b) &
\((s/8)\sum_j\Delta D_j+(J/2)\sum_j\Delta X_j\ge|\hbar|\arcsin(1-2\eta)\)
& Regular quantum joint states, unitary instrument, test valid for every
initial body state; the theorem's interpolating-state suprema \\
Theorem E & \(\lambda_0(ab/|\hbar|)\ge(1-2\eta)^2\) & Quantum state
concentrated in both position and momentum windows with probability at
least \(1-\eta\) each \\
\end{longtable}

Here \(s=F\tau^2/(2m)\), \(J=F\tau\), \(a,b\) are window half-widths,
and \(\lambda_0\) is the largest eigenvalue of the
time-and-band-limiting operator with the fifth-postulate note's
convention. For Theorem E's finite threshold take \(0<\eta<1/2\); at
\(\eta=0\) simultaneous compact concentration is impossible for finite
\(a,b\). A phase difference of \(\pi\) already gives perfect
discrimination in (7).

The geometry relates \(sJ=12K_\tau\). Converting Theorem C into a bound
on \(K_\tau\) additionally needs bounds on the allowed disturbances
\(\Delta D_j,\Delta X_j\) relative to \(J,s\). Converting Theorem E
needs a physical identification of the concentration widths \(a,b\) with
trajectory resolution. Choosing, for example, \(a=\alpha s\) and
\(b=\beta J\) would give \(K_\tau\ge|\hbar|c_*(\eta)/(12\alpha\beta)\),
with \(c_*(\eta)=\lambda_0^{-1}((1-2\eta)^2)\) and the positive
dimensionless apertures \(\alpha,\beta\) supplied. Geometry alone fixes
neither choice.

The phase experiment has a two-dimensional record space. Its
Fubini--Study distance from the reference is \(\delta/2\), and the
equal-prior error target requires that distance to be at least
\(\arcsin(1-2\eta)\). This is the same statistical-angle step used in
Theorem C's disturbance proof. Theorem E uses the overlap norm of two
spectral projections in an infinite-dimensional canonical state space.
Its concentration function and error dependence carry different
information. Their common right-hand scale comes from the supplied
canonical phase constant; the equality of the scale alone establishes no
equivalence of the three experimental tasks.

There is a stronger exact bridge within the restricted Gaussian mark
model. The \href{polygon-lift-phase.md}{polygon-lift note, Theorem 7}
(passage) uses the same functional
\(\mathcal K_\tau[f]=(2m)^{-1}\iint G_\tau(s,u)f(s)f(u)\,ds\,du\) for
the polygon phase and the optimal mark deflection. For constant force,
\(\mathcal K_\tau=K_\tau\); with mark trade-off \(\kappa=|\hbar|/2\),
its supremal squared deflection is \(d_{\rm mark}^2=2K_\tau/|\hbar|\).
That equality retains the Gaussian, uncorrelated error/recoil and
dense-protocol hypotheses. It leaves the general C and E bounds with
their own quantities and quantifiers.

Thus the two-path calculation gives the polygon theorem's existing phase
bound in the halved-functional representation, plus its explicit limit
prescriptions. It adds no independent exclusion of the common zero
branch allowed by \href{rotation-composition-universality.md}{round 2,
§4} (passage). Every \(\varepsilon>0\) supports (3)--(14); their
thresholds scale with that freely supplied value. At zero, the
exponential representation is singular, while classical mechanics still
supplies ordinary trajectories and, with state completeness, sharp
pointer records as in Theorem C(c). A singular coordinate for that
classical theory creates no inconsistency in the theory itself.

\hypertarget{consequence-for-state}{%
\subsection{6. Consequence for STATE}\label{consequence-for-state}}

Item 2 gains an exact two-path error law on the Galileo cubic action,
with the endpoint conversion, phase aliases, resources and limiting
prescriptions explicit. Fixed-resolution refinement erases this phase
record; classical-first averaging or calibrated attenuation gives the
precise contrasting limits. The result locates a conditional record
threshold at \(K_\tau=\tau\Delta E/12\) for the matched polygon
protocol, and at the corresponding calibrated action for a literal
inertial reunion. The remaining necessity step is to justify a common
positive phase/readout scale independently of the quantum premises and a
fixed experimental budget. Round 3 supplies the representation and its
exact obstruction, while that physical scale-selection problem remains
open.
\end{document}
